\label{sec:methodology}

We use a common finite-element notation across all benchmark families. The
continuous domain is $\Omega \subset \mathbb{R}^d$, the conforming trial space
is $V_h \subset V$, the state vector is $u_h \in V_h$, and the element-restricted
state on element $e$ is $u_e$. When a family is posed as an energy or reduced
objective minimization, we write its discrete functional as
$\mathcal{J}_h(u_h)$; when a family is posed as a sequence of equilibrium
solves or by a scalar equilibrium surrogate, we write the associated scalar
functional as $\Pi_h(u_h)$.
Hyperelasticity is the only family in which $J=\det F$ denotes the deformation
Jacobian, so we avoid the bare symbol $J$ for generic objectives elsewhere in
the paper.
For finite-element labels, $P_k(L_\ell)$ denotes degree-$k$ Lagrange elements
on mesh level $L_\ell$; problem-specific sections define the mesh hierarchy
behind those levels when it matters for interpretation.

At the implementation level, the common sparse-assembly task is therefore one
of two closely related forms:
\begin{equation}
  \mathcal{J}_h(u_h)
  = \sum_{e=1}^{n_{\mathrm{el}}} \Phi_e(u_e) - f_{\mathrm{ext}}^T u_h,
  \label{eq:generic-objective}
\end{equation}
or
\begin{equation}
  \Pi_h(u_h)
  = \sum_{e=1}^{n_{\mathrm{el}}} \Phi_e(u_e) - f_{\mathrm{ext}}^T u_h,
  \label{eq:generic-potential}
\end{equation}
where $\Phi_e$ may represent a scalar diffusion energy, a phase-field
potential, a finite-strain stored-energy density, a repository scalarized
plasticity surrogate, or the reduced design objective used in topology
optimization. The benchmark
sections specialize \cref{eq:generic-objective,eq:generic-potential} to each
family and define any additional symbols locally.

\subsection{Finite-element and derivative structure}

The main design choice in \repo{} is that problem definitions should remain
energy-centered whenever possible. This gives three derivative routes:
\begin{enumerate}
  \item \textbf{element AD:} differentiate the full scalar element energy
        $\Pi_e(u_e)$ with respect to the element DOF vector;
  \item \textbf{constitutive AD:} differentiate a quadrature-point energy
        density $\psi(\varepsilon_q)$, then assemble
        $r_e = \sum_q w_q B_q^T \sigma_q$ and
        $K_e = \sum_q w_q B_q^T C_q B_q$;
  \item \textbf{colored SFD:} recover sparse Hessian information through
        directionally probed gradients/HVPs and graph coloring.
\end{enumerate}

Element AD is the most literal route and is often the cleanest mathematically.
Constitutive AD preserves the same scalar-energy source of truth but changes
the differentiation boundary from the full element vector to the local strain
state. This is especially effective for high-order plasticity, where the full
element Hessian is often much more expensive than the $6\times6$
constitutive tangent. Colored SFD is used when exact second-order information
is too costly or unnecessary.

\subsection{Newton--Krylov globalization}

The nonlinear solves in the repository are based on Newton or quasi-Newton
updates, but they are not all globalized in the same way. Scalar benchmarks
primarily use line-search Newton. Hyperelasticity uses a trust-region method
with optional post-step line search. Some maintained experiments use a hybrid
trust-region/line-search algorithm in which a reduced trust-region model
provides the direction, while a bounded line search chooses the accepted step.
The line-search and trust-region ingredients follow standard globalization ideas
\citep[Chs.~2--4]{nocedal2006numerical}; the particular hybrid coupling used in
some repository experiments is an implementation-specific combination of those
textbook ingredients.

\begin{algorithm}[H]
\caption{Hybrid trust-region / line-search Newton}
\begin{algorithmic}[1]
\State Given $x_0$, trust radius $\Delta_0$, tolerances, and callbacks for
energy, gradient, and Hessian solves/matvecs
\For{$k=0,1,\dots$}
  \State Assemble $g_k = \nabla J(x_k)$
  \If{convergence criteria are satisfied}
    \State terminate
  \EndIf
  \State Compute a Newton-like direction from $H(x_k)p = -g_k$
  \If{trust region is enabled}
    \State Build a small reduced model in the span of Newton and gradient directions
    \State Solve the reduced trust subproblem with radius $\Delta_k$
    \State Optionally compare against a pure gradient trust step
    \State Restrict the line-search interval so the accepted step stays in the trust ball
  \Else
    \State Use the Newton direction directly
  \EndIf
  \State Perform one-dimensional search on the chosen direction
  \State Compute actual and predicted reduction
  \If{trust region is enabled}
    \State accept/reject using $\rho = \mathrm{ared}/\mathrm{pred}$ and update $\Delta_k$
  \Else
    \State accept only if the merit function decreases
  \EndIf
\EndFor
\end{algorithmic}
\end{algorithm}

For the maintained 3D plasticity path promoted in this paper, the chosen policy
is simpler: Newton with Armijo backtracking and a PMG-preconditioned
\code{fgmres} linear solve \citep[Sec.~3.1]{nocedal2006numerical}. This is the
maintained baseline workflow used for the flagship benchmark reported later in
\Cref{sec:num-plasticity3d}.

\begin{algorithm}[H]
\caption{Armijo line search used in maintained Newton solves}
\begin{algorithmic}[1]
\State Given state $x$, direction $p$, initial step $\alpha_0$
\State Evaluate $J(x)$ and $g = \nabla J(x)$
\For{$m=0,1,\dots,m_{\max}$}
  \State $\alpha \gets \alpha_0 \beta^m$
  \State $x_{\mathrm{trial}} \gets x + \alpha p$
  \If{$J(x_{\mathrm{trial}}) \le J(x) + c_1 \alpha g^T p$}
    \State accept $\alpha$
    \State \Return
  \EndIf
\EndFor
\State fall back to the smallest admissible step or to gradient-safe logic
\end{algorithmic}
\end{algorithm}

\subsection{Colored sparse finite differences}

Sparse derivative recovery is built around the observation that if columns that
do not interfere in the sparsity pattern are perturbed together, one gradient
or HVP evaluation can recover multiple columns of the Jacobian or Hessian
\citep{curtis1974sparsejacobian,coleman1984sparsehessian}. In this repository,
the local SFD path uses distance-2 coloring on the Hessian sparsity graph,
which is especially effective when combined with overlap-domain assembly and
owned-plus-ghost PETSc layouts.

\begin{algorithm}[H]
\caption{Distributed colored Hessian recovery}
\begin{algorithmic}[1]
\State Build or import local sparsity information for the free-DOF graph
\State Compute a distance-2 coloring of the Hessian pattern
\For{each color group $c$}
  \State Build the combined probe vector $v_c$
  \State Evaluate one gradient difference or HVP for $v_c$
  \State Scatter recovered entries into owned sparse storage
\EndFor
\State Assemble the global PETSc matrix from owned COO/CSR entries
\end{algorithmic}
\end{algorithm}

\subsection{Constitutive autodiff for 3D plasticity}

The promoted \plasticitythree{} result in this paper uses constitutive AD.
Rather than applying second-order AD to the full $P4$ element vector, we write
one scalar Mohr--Coulomb energy density per quadrature point,
\begin{equation}
  \Phi_q(\varepsilon_q,\mathrm{material}_q),
\end{equation}
differentiate it with respect to the six engineering-strain components, and
assemble the element tangent from the resulting stress and constitutive tangent.
This keeps the scalar-energy contract intact while reducing the cost of
branchwise second-order information on high-order tetrahedra.

\begin{algorithm}[H]
\caption{Constitutive AD assembly for high-order plasticity}
\begin{algorithmic}[1]
\For{each local element $e$}
  \For{each quadrature point $q$}
    \State Evaluate strain $\varepsilon_q = B_q u_e$
    \State Evaluate scalar energy density $\Phi_q(\varepsilon_q)$
    \State Obtain stress $\sigma_q = \partial \Phi_q / \partial \varepsilon_q$
    \State Obtain constitutive tangent $C_q = \partial^2 \Phi_q / \partial \varepsilon_q^2$
  \EndFor
  \State Assemble $r_e = \sum_q w_q B_q^T \sigma_q$
  \State Assemble $K_e = \sum_q w_q B_q^T C_q B_q$
\EndFor
\end{algorithmic}
\end{algorithm}

The key methodological point is that automatic differentiation is applied
directly to the implemented scalar branch potential. The code does not switch
to a hand-derived constitutive tangent. Away from branch switching sets this
yields branchwise exact derivatives of the implemented surrogate; at switching
sets the model remains nonsmooth, so the resulting tangents should be read as
branchwise surrogate derivatives rather than globally smooth continuum
derivatives.
